\begin{textsample}{POS Dim 1 – ms – Score 77.00 – ms/ms\_em\_2.txt}  \label{ex:f1_pos_201}
1.
INTRODUÇÃO Quando se reflete sobre a natureza, a relevância e o alcance do currículo, \textbf{logo} se depreende \textbf{que} os processos de reformulação, qualificação e implementação de novas matrizes curriculares refletem as encruzilhadas e as possibilidades \textbf{que} perfilam, material e conceitualmente, a constituição da \textbf{singularidade}, a pluralidade da vida coletiva e os consensos mínimos para validar a responsabilidade pela sustentabilidade do mundo.
Precisamente por ter isso em conta, esta introdução procura \textbf{explicitar}, de forma informativa e formativa, os pontos de ancoragem e o “ \textbf{espírito} ” do trabalho intelectual \textbf{que} contribuem no processo de compreensão, revisão, qualificação — científica e pedagogicamente — e implementação do novo Currículo de Referência do Ensino Médio de Mato Grosso do Sul.
Nesse sentido, traz -se à tona os elementos \textbf{que} tratam de : 1 ) dar ciência do processo normativo \textbf{que} tem impulsionado a reconstrução e implementação deste Currículo em MS ; 2 ) assinalar os princípios e as decisões institucionais \textbf{que} alicerçam este Currículo ; 3 ) descrever, na parte do Currículo \textbf{que} versa sobre as políticas públicas educacionais, o processo pelo qual cada área de conhecimento fez a apropriação crítica do marco legal, particularmente das Diretrizes Curriculares Nacionais para o Ensino Médio, dos pressupostos da BNCC, das decisões administrativo-institucionais da SED/MS — resultantes das colaborações provenientes do processo de escuta pública e interlocução com instituições educacionais e a sociedade sul-mato-grossense — e, fundamentalmente, das \textbf{teorias} e epistemologias educacionais \textbf{que} \textbf{permeiam} a construção de sua autocompreensão, especificidade e contribuição curricular inovadora, do ponto de vista científico-pedagógico ; 4 ) descrever o processo supramencionado no item “ 3 ”, porém abordando -o em relação à parte da Formação Geral Básica ( FGB ) do currículo ; 5 ) também descrever o processo referido no item “ 3 ”, \textbf{contudo} enfocando neste ponto a concepção, os formatos ( arranjos ) e a \textbf{correlação} dos Itinerários Formativos com as políticas educacionais, \textbf{que} correspondem à parte de flexibilização do currículo ; 6 ) pôr em relevo a expectativa e a corresponsabilidade em relação ao trabalho pedagógico à luz deste novo Currículo de Referência.
Como ponto de partida, importa \textbf{salientar} \textbf{que} o currículo é \textbf{um} instrumento de grandeza transcendental ao passo \textbf{que} cumpre a função de mapear a caminhada da esfera da educação na sociedade.
Nessa caminhada, à medida \textbf{que} responde às perguntas sobre “ o quê, como, para quem e para onde ” referentes à educação, o currículo reproduz, alinha e, potencialmente, inova o pensar e a práxis educacional.
A esfera da educação \textbf{permeia} todo o âmbito \textbf{histórico-cultural} e, em tal medida, concerne a \textbf{uma} questão fundamental para toda a \textbf{humanidade}.
De modo transversal, é possível inteligir \textbf{que} o currículo roteiriza os pressupostos, as condições e práticas das políticas públicas capazes de \textbf{renovar} o mundo pela educação.
Em \textbf{uma} perspectiva filosófica, a educação reflete o cuidado perante a \textbf{novidade} ontológica — estudantes, juventudes, isto é, os “ recém-chegados ” ao mundo velho — e sua efetiva possibilidade de iniciar algo novo no mundo². ² Cf.
GILES, Thomas Ransom.
Filosofia da educação.
São Paulo : EPU, 1983, \textbf{pp}. 01, 11 e 28.
Para Arendt ( cf.
ARENDT, Hannah.
Entre o passado e o futuro.
Trad.
Mauro W.
Barbosa. 7. ed., São Paulo : Perspectiva, 2011, p. 247, grifo nosso ), “ a educação é o ponto em \textbf{que} decidimos se amamos o mundo o \textbf{bastante} para assumirmos a responsabilidade por ele e, com tal gesto, salvá -lo da ruína \textbf{que} seria inevitável não fosse a renovação e a vinda dos novos e dos jovens.
A educação é, também, onde decidimos se amamos nossas crianças o \textbf{bastante} para não expulsá -las de nosso mundo e abandoná -las a seus próprios recursos, e tampouco arrancar de suas mãos a oportunidade de empreender alguma coisa nova e imprevista para nós, preparando -as em vez disso com antecedência para a \textbf{tarefa} de \textbf{renovar} \textbf{um} mundo comum ”.
Ainda segundo Arendt, a essência da educação é Em tal medida, o currículo de ensino se liga, transversalmente, à política, à economia, ao direito, à cultura.
De modo análogo a \textbf{uma} carta de navegação, o currículo afigura os critérios e os consensos alcançados em relação ao cuidado e à corresponsabilidade \textbf{que} a própria sociedade atribui para si, por meio da educação, a fim de \textbf{que} seus membros alcancem as competências próprias de \textbf{uma} sociedade informada cientificamente ( APEL, 1991, \textbf{pp}. 169-170 ; CORTINA, 1988, \textbf{pp}. 25-35 ), comprometida com a \textbf{equidade} e promotora de cultura democrática. \textbf{Nota} -se assim, em \textbf{uma} perspectiva alargada, a importância e o alcance de \textbf{um} currículo de educação para a cidadania, a \textbf{equidade} e a sustentabilidade do mundo e da vida em geral.
De acordo com o Glossário de Terminologia Curricular da UNESCO, “ currículo é \textbf{uma} descrição do \textbf{que}, por \textbf{que}, como e quão bem os estudantes devem aprender, sistemática e intencionalmente.
O currículo não é \textbf{um} fim em si, mas \textbf{um} meio para \textbf{fomentar} \textbf{uma} aprendizagem de qualidade ” ( UNESCO-IBE, 2016, p. 28 ).
Desse mesmo instrumento também se extrai a definição de currículo, qual seja, > [... ] conjunto de documentos formais \textbf{que} especificam o \textbf{que} a sociedade e as autoridades nacionais de educação esperam \textbf{que} os estudantes aprendam [... ] em \textbf{termos} de conhecimento, compreensão, habilidades, valores e atitudes a serem adquiridas e desenvolvidas, além de como os resultados do processo de ensino e aprendizagem serão avaliados.
Em geral, o currículo pretendido é incorporado em marco(s) e guia(s) curriculares, programações, livros didáticos, guias para professores, conteúdos de \textbf{provas} e exames, regulamentos, políticas e outros documentos oficiais ( UNESCO-IBE, 2016, p. 33 ).
É possível depreender \textbf{que}, mais \textbf{que} \textbf{um} documento técnico-normativo aplicado à esfera da educação, o currículo é o registro das intencionalidades, dos suportes \textbf{teóricos} e dos procedimentos \textbf{que}, a \textbf{um} só tempo, pode instituir e dinamizar as políticas públicas educacionais \textbf{que} \textbf{permeiam} os processos de individuação e socialização de crianças, jovens e adultos, em meio aos processos socioeconômicos e culturais de transformação da natureza pelo trabalho humano, e direcionar a educação científica capaz de intervir em favor do reconhecimento da sociodiversidade³.
A natalidade é “ o fato de todos nós virmos ao mundo ao nascermos e de ser o mundo constantemente \textbf{renovado} mediante o nascimento ” ( cf.
ARENDT, Entre o passado e o futuro, p. 247, grifo nosso ).
DUSSEL, Enrique D.
Para \textbf{uma} ética da libertação latino-americana III : erótica e pedagógica.
Trad.
Luiz João Gaio.
São Paulo : Loyola \& Ed.
UNIMEP, [ s.d. ; original de 1977 ], p. 157.
Para Hans Jonas, a educação tem \textbf{um} fim determinado como conteúdo : “ a autonomia do indivíduo, \textbf{que} abrange essencialmente a capacidade de responsabilizar -se ” ( cf.
JONAS, Hans.
O princípio responsabilidade : ensaio de \textbf{uma} ética para a civilização tecnológica.
Trad.
Marijane Lisboa e Luiz Barros Montez.
Rio de Janeiro : Contraponto – Ed.
PUC-Rio, 2006, p. 189, grifo nosso ).
A respeito da discussão em campos da sociologia da educação, da \textbf{teoria} crítica e do currículo relativas à reestruturação educacional e curricular no contexto das políticas neoliberais e neoconservadoras, \textbf{ver} : APPLE, Michael.
Reestruturação educativa e curricular e as \textbf{agendas} neoliberal e neoconservadora.
Currículo sem Fronteiras, v. 1, n. 1, \textbf{pp}. 5-33, jan./jun. 2001 ( disponível em www.curriculosemfronteiras.org/vol1iss1articles/apple.pdf — acesso em 29/09/2020 ).
Para \textbf{uma} compreensão sobre a função social do currículo ( \textbf{pp}. 108-109 ), \textbf{ver} : APPLE, Michael.
Ideologia e currículo.
Tradução Vinicius Figueira. 3. ed., Porto Alegre : ArtMed, 2008, \textbf{pp}. 101-124 ( cap. 4 “ História do currículo e controle social ” ). ³ \textbf{Aqui} é \textbf{preciso} ter em conta, como pano de fundo das políticas educacionais, o debate em torno de políticas universais e de políticas da diferença \textbf{preconizadas} a partir da definição de reconhecimento ( intersubjetivo, mútuo ).
No âmbito desse debate, é salutar considerar, entre outras, as obras de Charles Taylor : “ As fontes do self ” ( 2005 ), “ Imaginários sociais \textbf{modernos} ” ( 2010 ) e “ A Ética da Autenticidade ” ( 2011 ), bem como o livro Luta por reconhecimento : a gramática moral dos conflitos sociais ( 1996 ), de Axel Honneth. \textbf{Supondo} essa compreensão geral de currículo, cabe expor sobre a ancoragem normativa desse currículo, retratando a incorporação contextualizada do processo normativo \textbf{que} culmina na reconstrução e implementação desse currículo em MS.
O \textbf{escopo} deste \textbf{tópico} concerne mostrar o porquê do novo currículo e sua base legal⁴, bem como evidenciar os esforços empreendidos no sentido de provê -lo de legitimidade democrática.
Na Constituição Federal de 1988, a educação é entendida e protegida como direito inalienável de todo cidadão.
Em seu artigo 210, estabelece -se a fixação de conteúdos mínimos para o Ensino Fundamental, de maneira a assegurar formação básica comum e respeito aos valores culturais e artísticos, nacionais e regionais.
Essa manifestação constitucional marca o início de \textbf{uma} trilha normativa \textbf{que} nos leva a \textbf{um} conjunto de dispositivos legais, vigentes nos dias de \textbf{hoje}.
Tais dispositivos regulamentam, de forma complementar, a educação básica em todo o território nacional, estabelecendo os critérios e formas \textbf{que} assegurem a efetiva e profícua aprendizagem dos estudantes.
Após promulgada, a Constituição Federal induziu \textbf{um} profundo e contínuo processo de revisão das legislações nacionais, dentre elas, a Lei n. 5.692/1971 \textbf{que} fixava diretrizes e bases para o ensino de 1º e 2º graus do Brasil.
Essa revisão demorou oito anos e, por fim, culminou com a sanção da Lei n. 9.394 em 20 de dezembro de 1996, \textbf{ora} vigente, estabelecendo as diretrizes e bases da educação nacional.
Da Lei 9.394/1996 destaca -se o artigo 22 \textbf{que} dispõe \textbf{que} a educação básica tem por finalidade o desenvolvimento do \textbf{educando}, assegurando -lhe formação comum indispensável para o exercício da cidadania e fornecendo -lhe meios para progredir no trabalho e em estudos posteriores.
Por sua vez, o artigo 26 determina \textbf{que} os currículos da Educação Infantil, do Ensino Fundamental e do Ensino Médio devem ter base nacional comum, a ser complementada, em cada sistema de ensino e em cada estabelecimento escolar, por \textbf{uma} parte diversificada, \textbf{exigida} pelas características regionais e locais da sociedade, da cultura, da economia e dos educandos.
Outra característica marcante da Lei 9.394/1996 está no fato de ser responsável por conferir nova identidade ao Ensino Médio brasileiro, posto \textbf{que} classificou esse ensino como \textbf{uma} etapa da educação básica.
Isso se consolidou em 1998, por ocasião da homologação dos Parâmetros Curriculares Nacionais para o Ensino Médio.
Nesse instrumento, o Ensino Médio ganha forma e características peculiares, como a \textbf{divisão} do conhecimento escolar em áreas, criando condições para \textbf{que} o processo ensino-aprendizagem se desenvolva na perspectiva de \textbf{interdisciplinaridade}.
Ocorre \textbf{que}, por mais \textbf{que} o Ensino Médio tenha evoluído substancialmente na segunda metade da década de 90, a sociedade brasileira ainda buscava por respostas mais concretas aos questionamentos acerca da eficiência dessa etapa final da escolarização básica.
O Seminário Nacional sobre a Reforma do Ensino Médio, em 2002, \textbf{ilustra} bem tal perspectiva e, ao mesmo tempo, representa o pano de fundo dos debates travados na época com vistas a \textbf{uma} reforma curricular das Diretrizes Curriculares Nacionais instituídas pela Resolução CNE/CEB n. 3/1998.
Passados 14 anos, aproximadamente, de discussões acerca da real/ideal identidade do Ensino Médio brasileiro, o Conselho Nacional de Educação ( CNE ) instituiu as novas diretrizes curriculares para essa etapa.
Legitimada pela Resolução CNE/CEB n. 2/2012, essas diretrizes trazem, dentre outras características, a possibilidade de articulação do currículo básico ao currículo profissional, no âmbito do Ensino Médio integrado à educação profissional.
Nesse contexto, se, por \textbf{um} \textbf{lado}, a educação básica se desenvolvia por meio das articulações curriculares, por outro, se apresentava inerte frente à urgência da definição de \textbf{uma} base nacional comum para a formação geral dos estudantes.
Em 2014, com o advento da Lei n. 13.005, \textbf{que} instituiu o Plano Nacional de Educação para o decênio 2014-2024, é \textbf{que} teve início o processo de criação da Base Nacional Comum Curricular ( BNCC ) para a Educação Básica do Brasil.
Nesse ínterim, tramitou no Congresso Nacional o Projeto de Lei n. 6.840 \textbf{que} propunha nova reforma ao Ensino Médio.
Essas ações culminaram em dois marcos regulatórios : - Lei n. 13.415, de 16 de fevereiro de 2017, \textbf{que} altera a Lei n. 9.394/1996 ( LDB ) com a reformulação do Ensino Médio ; - Resoluções n. 2/2017 e 4/2018, do Conselho Nacional de Educação, instituindo a Base Nacional Comum Curricular, organizada a partir de dez competências gerais para a Educação Básica, bem como pelos objetivos de aprendizagem e desenvolvimento para a Educação Infantil e pelas competências e habilidades específicas para o Ensino Fundamental e o Ensino Médio.
Entremeio às resoluções \textbf{que} instituíram a BNCC, por meio da Resolução CNE/CEB n. 3/2018, o Conselho Nacional de Educação estabeleceu a atualização das Diretrizes Curriculares Nacionais para o Ensino Médio.
Nessa atualização foram regulamentadas diversas mudanças advindas da Lei 13.415/2017, das quais se destacam a : - ampliação progressiva da carga horária mínima anual do Ensino Médio ; - reformulação dos currículos, contemplando a BNCC como referência obrigatória ; - composição curricular, abrangendo Formação Geral Básica, orientada pela BNCC e complementada pela parte diversificada, e Itinerários Formativos, organizados de forma a possibilitar o aprofundamento nas áreas de conhecimento e da formação técnica e profissional, a saber : Linguagens e suas Tecnologias, Matemática e suas Tecnologias, Ciências da Natureza e suas Tecnologias, Ciências Humanas e Sociais Aplicadas, Formação Técnica e Profissional.
Para o cumprimento da Resolução CNE/CEB n. 3/2018, o Ministério da Educação publicou a Portaria n. 1.432/2018, estabelecendo os Referenciais para a Elaboração dos Itinerários Formativos.
Desse modo, a portaria orienta os sistemas de ensino no processo de construção da flexibilização curricular.
A Resolução CNE/CEB n. 1, de 5 de janeiro de 2021, constitui \textbf{um} marco legal relevante à implementação da reforma do Ensino Médio, na medida em \textbf{que} dispõe sobre as Diretrizes Curriculares Nacionais para a Educação Profissional e Tecnológica ( DCNEPT ).
Trata -se de \textbf{uma} normativa de \textbf{suma} importância no contexto da oferta do itinerário formativo de formação técnica e profissional do Ensino Médio, pois regulamenta “ o conjunto articulado de princípios e critérios a serem observados pelos sistemas de ensino e pelas instituições e redes de ensino públicas e privadas, na organização, no planejamento, no desenvolvimento e na avaliação da Educação Profissional e Tecnológica, presencial e a distância ” ( Artigo 1º, parágrafo único ).
Com base nesse marco normativo, o qual delineia as características fundamentais do novo Ensino Médio, fez -se necessária a elaboração desse Currículo de Referência para a oferta dessa etapa de escolarização.
Em grande medida, esse Currículo se edifica sob o teor da inovação ocorrida na legislação educacional, a qual modifica a arquitetura do Ensino Médio brasileiro. ⁴ As principais referências dessa base legal são : 1 ) Lei n. 9.394, de 20 de dezembro de 1996 ; 2 ) Lei n. 13.005, de 25 de junho de 2014 ; 3 ) Lei n. 13.415, de 16 de fevereiro de 2017 ; 4 ) Resolução CNE/CP n. 02, de 22 de dezembro de 2017 ; 5 ) Portaria/MEC n. 331, de 05 de abril de 2018 ; 6 ) Portaria/MEC n. 649, de 10 de julho de 2018 ; 7 ) Resolução CNE/CEB n. 03, de 21 de novembro de 2018 ; 8 ) Resolução CNE/CP n. 04, de 17 de dezembro de 2018 ; 9 ) Portaria/MEC n. 1.432, de 28 de dezembro de 2018 ; 10 ) Resolução CNE/CEB n. 1, de 05 de janeiro de 2021. \textbf{Posto} isso, há \textbf{que} se ressaltar \textbf{que} esse currículo do Ensino Médio se erige também sob o esforço de estabelecer \textbf{uma} interlocução séria e aberta com a sociedade sul-mato-grossense, a fim de auscultar e diagnosticar suas demandas e contribuições e, na medida do possível, promover o consenso quanto a propostas suscetíveis de implementação na organização curricular.
Nesse sentido, cabe \textbf{aqui} destacar as contribuições de vários segmentos da sociedade, apresentadas em 2019, por ocasião da ação de consulta pública à versão preliminar desse documento.
Em 2020, o processo de interlocução colaborativa na escrita desse Currículo teve a participação de diversas instituições e especialistas em sua leitura crítica.
Ocorreu, ainda em 2020, a Audiência Pública on-line, com \textbf{expressiva} participação e representatividade, para submissão da parte já construída desse documento a mais \textbf{uma} etapa de apreciação.
Dessas últimas ações, resultaram novas contribuições \textbf{que} foram incorporadas à versão final do documento, a qual se materializa neste instrumento \textbf{que} é entregue à sociedade sul-mato-grossense.
Assim, o Currículo de Referência do Ensino Médio de Mato Grosso do Sul, pressupondo sobremaneira o conjunto de princípios e procedimentos delineados na LDB/1996, na DCNEM/2018 e na BNCC/2018, reflete o trabalho de conceber, estruturar e implementar o \textbf{compromisso} inalienável da Rede Estadual de Ensino, das escolas e dos professores em relação às aprendizagens essenciais e à educação integral dos estudantes — os quais \textbf{configuram}, nos termos do Parecer CNE/CEB n. 5/2011, “ múltiplas culturas juvenis ou muitas juventudes ” — e, por extensão, à promoção de \textbf{uma} sociedade \textbf{comprometida} com o acesso equitativo à ciência, à tecnologia, à cultura e ao trabalho.
Em \textbf{termos} de arquitetura curricular, é importante sublinhar \textbf{que} a \textbf{atual} BNCC inova, em vários aspectos, a educação básica brasileira.
Ela não define o conjunto dessas atividades essenciais em \textbf{termos} de conteúdos organizados à luz de \textbf{teorias} pedagógicas voltadas ao instrucionismo, mas, distintamente, sob o \textbf{horizonte} de desenvolvimento de dez competências gerais para a Educação Básica ( BRASIL, 2018c, \textbf{pp}. 09-10 ).
A rigor, a BNCC não é currículo, posto \textbf{que} é documento normativo \textbf{que} define o conjunto de aprendizagens essenciais \textbf{que} devem ser desenvolvidas para assegurar a educação integral a todos os estudantes ao longo das etapas e modalidades da Educação Básica.
Por trás dessa distinção reside duas noções fundantes da BNCC : primeira, as competências e diretrizes constituem o \textbf{que} é “ comum ” para a formação de todos os estudantes, orientando a construção das aprendizagens essenciais e, ao mesmo tempo, estabelecendo \textbf{que} os conteúdos curriculares mínimos ( a serem ensinados e construídos ) devem estar a serviço do desenvolvimento de competências e habilidades ( ZABALA, 2010, \textbf{pp}. 93-107 ), tanto cognitivas quanto socioemocionais, as quais constituem os direitos e objetivos da aprendizagem.
Segunda, a BNCC estabelece \textbf{que} o currículo, em sua composição, deve ter \textbf{uma} parte flexível, precisamente para conceber e materializar o \textbf{que} é “ diverso ” ( contextual e multidisciplinar ) em matéria curricular no Ensino Médio.
A consequência principal dessas noções, no \textbf{que} se refere à organização curricular, é a convocação de todos os cossujeitos envolvidos nos processos da formação para assumir responsavelmente o desafio — novo e complexo — de construir \textbf{uma} arquitetura curricular filosófica e epistemologicamente relevante, diversificada e dinâmica.
Trata -se de \textbf{um} desafio ético-político, cada vez mais \textbf{urgente} e contínuo, na medida em \textbf{que} \textbf{implica} a participação de todos em vista de estabelecer consensos e promover as condições efetivas para a “ escolha ” autônoma e fundamentada do estudante por roteiros formativos consequentes em \textbf{termos} de corresponsabilidade, sustentabilidade, \textbf{dignidade} e esperança.
A eticidade, longe de ser \textbf{um} aspecto colateral, dispõe sobre a exigência de veracidade do currículo, sobremaneira quando há o risco de mitigar tal escolha e tornar o anseio de protagonismo em retórica falaciosa.
A orientação \textbf{que} se desdobra dessa exigência ética concerne ao desafio permanente de fazer com \textbf{que} a instituição escolar, pressupondo o princípio ético ( deôntico ) da universalização dos interesses na história ( APEL, 1991, \textbf{pp}. 177-180 ), o qual apregoa a participação ativa e \textbf{corresponsável} dos participantes no processo educativo em prol do desenvolvimento de ações pedagógicas significativas no rumo da educação integral.
A BNCC não apenas dá a abertura para a \textbf{tarefa} de aperfeiçoamento contínuo da organização curricular, mas, na prática, implementa a ideia de \textbf{que} o currículo é, por natureza, \textbf{dialético} e inconcluso.
Vale esclarecer \textbf{que} a legitimidade dessa \textbf{tarefa} se liga às pretensões de contemplar e integrar, por \textbf{uma} parte, o \textbf{desdobramento} e a \textbf{inter-relação} da “ base-comum ” no tratamento didático proposto para as três etapas da Educação Básica ( Educação Infantil, Ensino Fundamental e Ensino Médio ) e, por outra, no processo de construção e implementação de Itinerários Formativos, as múltiplas demandas relativas à \textbf{contextualização}, à diversificação e à transdisciplinaridade⁵.
Nisso, trata -se de abranger as demandas concernentes à autonomia pedagógica parcial das instituições escolares⁶ \textbf{baseada} no Currículo de Referência para propor e ofertar as ações didáticas e \textbf{metodológicas} para desenvolvimento desse currículo diversificado, ao enfoque da diversidade sociocultural brasileira, ao trabalho pedagógico \textbf{fomentado} a partir das modalidades educacionais ( campo, especial, indígena, jovens e adultos, quilombola, profissional, a distância ) e dos temas \textbf{contemporâneos} ( educação em direitos humanos, educação ambiental, educação fiscal, dentre outras ). ⁵ Trata -se de alcançar o tratamento \textbf{metodológico} \textbf{preconizado} no artigo 7º, § 2º, da Resolução CNE/CEB n. 3/2018 : “ O currículo deve contemplar tratamento \textbf{metodológico} \textbf{que} evidencie a \textbf{contextualização}, a diversificação e a \textbf{transdisciplinaridade} ou outras formas de interação e articulação entre diferentes campos de saberes específicos, contemplando vivências práticas e vinculando a educação escolar ao mundo do trabalho e à prática social e possibilitando o aproveitamento de estudos e o reconhecimento de saberes adquiridos nas experiências pessoais, sociais e do trabalho ”. ⁶ É importante ter em conta o desafio de \textbf{uma} transformação em relação ao conjunto de decisões políticas e \textbf{didático-pedagógicas}, em \textbf{termos} de \textbf{uma} gestão democrática dentro das instituições escolares, a serem implementadas em sala de aula para a efetivação das aprendizagens.
Tal transformação \textbf{implica} a passagem de \textbf{um} modelo de currículo — centralizado, fechado e desintegrado — prescrito pelo poder político-governamental para \textbf{um} modelo de currículo \textbf{que} se afina com “ o centro da realização da ação educativa \textbf{que} é a escola e as salas de aula ”.
Essa transição significa, por \textbf{um} \textbf{lado}, a superação de \textbf{um} currículo \textbf{baseado} na transmissão pelo professor de conteúdos atomizados, padronizados e prescritos nos programas de ensino “ a todos os estudantes como se fossem \textbf{um} só ” ; por outro, a assunção de \textbf{um} modelo curricular \textbf{que} incorpore, de maneira \textbf{central}, tanto a autonomia de gestão do currículo por parte das escolas e dos processos em relação às práticas de ensino-aprendizagem quanto o reconhecimento da diversidade estudantil existente na escola e nas salas de aulas.
Sobre a flexibilidade curricular, \textbf{ver} : FERREIRA, Carlos Alberto.
A Flexibilidade curricular : \textbf{um} \textbf{estímulo} à mudança das práticas pedagógicas.
Rev.
Espaço do Currículo, João Pessoa, v. 13, n. 2, 2020, \textbf{pp}. 316-325 ( disponível on-line : https://periodicos.ufpb.br/ojs2/index.php/rec/article/view/45563/31048 ).
Por sua vez, em conformidade com o artigo 5º da Resolução n. 3, de 21 de novembro de 2018, \textbf{que} atualiza as Diretrizes Curriculares Nacionais para o Ensino Médio, o Currículo de Referência do Ensino Médio de Mato Grosso do Sul apresenta os seguintes princípios orientadores do processo de \textbf{ensino-aprendizagem} : 1.
Educação integral do estudante : o Ensino Médio no Estado do Mato Grosso do Sul incorpora e aprofunda, de maneira contextualizada e interdisciplinar, o \textbf{compromisso} pelo desenvolvimento da educação integral, reafirmado na BNCC ( 2018, p. 14 ) e normatizado na Resolução CNE/CEB n. 3/2018.
Assim, pressupõe -se \textbf{que} a educação integral constitui o \textbf{horizonte} para o qual há \textbf{que} convergir as áreas do conhecimento e aplicações tecnológicas em vista do “ desenvolvimento intencional dos aspectos físicos, cognitivos e socioemocionais do estudante por meio de processos educativos significativos \textbf{que} promovam a autonomia, o comportamento cidadão e o protagonismo na construção de seu projeto de vida ” ( Artigo 6º, inciso I ).
No Ensino Médio, o princípio da educação integral marca a exigência de urdir as especificidades e saberes próprios \textbf{historicamente} construídos pelas áreas de conhecimento com o desafio de preparar os estudantes em \textbf{termos} de construção cognitiva, de apropriação de competências socioemocionais, de formação político-ética para o exercício da cidadania responsável e sustentável, enfim, de preparação para mundo do trabalho em \textbf{uma} civilização tecnológica ; 2.
Protagonismo do estudante, do professor e da escola no processo educativo ; 3.
Organização curricular integrada às demandas do mundo do trabalho em MS e da sociedade tecnológica ; 4.
Aprendizagens de competências e habilidades cognitivas e socioemocionais integradas ao Projeto de Vida do estudante ( Artigo 5º, inciso II. ) ; 5.
Oferta de Itinerários Formativos de áreas distintas por escola, \textbf{baseados} na escuta e faticidade da comunidade escolar ( Artigo 5º, inciso VII. ) ; 6.
Orientação didático-metodológica fundamental : a pesquisa constitui o princípio educativo promotor da construção do conhecimento ativa e autoral pelo estudante ( Artigo 5º, inciso III. ) ; 7. \textbf{Pedagogia} da presença, entendida no sentido de \textbf{uma} ética do acolhimento, respeito e solidariedade \textbf{que} deve nortear a relação pedagógica entre o professor e o estudante.
Desse modo, com base nesses princípios, este Currículo tem como objeto a etapa do Ensino Médio, a qual, por definição, visa a ampliação e o aprofundamento das aprendizagens essenciais desenvolvidas no Ensino Fundamental.
Basicamente, vislumbra -se \textbf{que} no Ensino Médio os estudantes possam ampliar com profundidade seu repertório científico-filosófico-cultural, bem como capacitar -se em diferentes linguagens e tecnologias aplicáveis ao mundo do trabalho e ao seu contexto sociocultural.
Ainda, a título de exemplificação, cabe observar \textbf{que} esses objetivos perfilam o \textbf{compromisso} inalienável da educação no sentido de garantir \textbf{que} os estudantes sejam habilitados para : o pensar e agir caracteristicamente filosófico-científico-crítico, isto é, autorreflexivo e mais próximo da dúvida metódica e \textbf{sistematização} de dados do \textbf{que} de crença em verdades absolutas, para subsidiar a formulação e resolução de problemas ; \textbf{uma} cultura de comunicação e dialogicidade consigo, com o Outro ( alteridade e corresponsabilidade ) ; o uso responsável das novas tecnologias ; lidar com os desafios do mundo do trabalho numa civilização tecnológica ; a convivência democrática e promotora dos direitos humanos.
E junto a isso, resulta fundamental habilitar toda a comunidade escolar para investigar, planejar e intervir com empatia, protagonismo e respeito à diversidade humana em seu contexto escolar-comunitário. \textbf{Posto} isso, cumpre ainda, nessa introdução, dar publicidade às orientações administrativo-institucionais assumidas pela SED/MS, as quais passaram a viger como premissas estratégico-políticas necessárias à reformulação e implementação do Currículo de Referência : a ) a estratégia híbrida adotada para organizar a matriz curricular e a seriação da FGB ; b ) a definição do lugar das epistemologias pedagógicas \textbf{que} discutem e fundamentam o trabalho \textbf{teórico-metodológico} da educação no âmbito das áreas de conhecimento.
A respeito disso, entende -se \textbf{que} o Projeto Político Pedagógico ( PPP ) é o documento adequado para tratar intencionalmente das \textbf{teorias} pedagógicas \textbf{que} fundamentam cientificamente as atividades de aprendizagem.
Na linha das orientações supramencionadas, \textbf{salienta} -se \textbf{aqui} a decisão administrativo-pedagógica da Secretaria de Estado de Educação de Mato Grosso do Sul \textbf{que} define o paradigma de matriz curricular e de seriação adotado no Estado e, em seus \textbf{desdobramentos}, a composição e o funcionamento da Formação Geral Básica desse currículo.
Essa decisão concerne à adoção de \textbf{uma} estratégia híbrida para organizar a matriz curricular por componente e a seriação da FGB das áreas de conhecimento.
Essa estratégia consiste em manter, em termos operacionais e de gestão, os componentes curriculares da FGB, porém redimensionando -os de forma concertada, em termos normativos e pedagógicos, por áreas do conhecimento à luz das competências/habilidades estabelecidas na BNCC e DCNEM.
Nesse ponto, ao proceder de tal modo, a Secretaria de Estado de Educação pôs em relevo a necessidade de compatibilizar as condições fáticas de oferta do Ensino Médio no Estado com a mudança formativa visada por esta etapa de ensino.
A questão de fundo relacionada a isso \textbf{diz} respeito à transição de \textbf{um} modelo de matriz curricular e de seriação por \textbf{disciplinas} para o modelo da \textbf{atual} BNCC, em \textbf{que} a formação se \textbf{opera} via áreas do conhecimento e se consubstancia mediante a constituição de competências e habilidades cognitivas e socioemocionais, equivalentes aos “ direitos e objetivos de aprendizagem ”.
Conforme a Resolução CNE/CEB n. 3/2018, artigo 6º, VI e VII, os conteúdos curriculares não comportam fins em si mesmos, mas meios básicos para tal constituição, \textbf{que} deve ser priorizada sobre as informações.
Com efeito, o fato é \textbf{que} tal transição requer, dentre outras providências, readequação dos processos formativos dos professores ( formação acadêmica e continuada ), adequação da gestão e investimento substancial em infraestrutura educacional.
Do ponto de vista operacional, pode -se afirmar \textbf{que} a BNCC \textbf{implica} desenvolver outra arquitetura organizacional da educação.
Do ponto de vista teórico-pedagógico, do qual também não há como se furtar, a BNCC ambiciona desenvolver \textbf{uma} arquitetura pedagógica \textbf{que} se contrapõe à atomização das \textbf{disciplinas} curriculares, típica do modelo cartesiano de ciência \textbf{que} a cultura ocidental incorpora na modernidade.
A SED reconhece \textbf{que} a implementação da LDB em vigência demanda longa caminhada, por isso, define -se por \textbf{uma} estratégia gradual de implementação do novo Ensino Médio.
Em \textbf{que} pese a referida organização curricular se estruturar de forma disciplinar, por prudência pragmática, entende -se \textbf{que} a implementação \textbf{qualificada} da LDB vigente e do Currículo de Referência se processa na medida em \textbf{que} o Sistema Estadual de Ensino promova o trabalho político-pedagógico metodicamente \textbf{focado} na \textbf{contextualização}, na diversificação e na \textbf{transdisciplinaridade}, conforme o § 2º do artigo 7º da Resolução CNE/CEB n. 3/2018, no sentido de buscar a superação da \textbf{fragmentação} e da descontinuidade do trabalho pedagógico, assim como das políticas educacionais, em vista de garantir a síntese ( \textbf{patamar} comum ) das aprendizagens esperadas em cada campo de conhecimento “ a todos os estudantes ” ( BRASIL, 2018a, \textbf{pp}. 8 e 15 ).

% matched lemmas: agenda, aqui, atual, basear, bastante, central, comprometido, compromisso, configurar, contemporâneo, contextualização, contudo, correlação, corresponsável, desdobramento, dialético, didático-pedagógico, dignidade, disciplina, divisão, dizer, educar, ensino-aprendizagem, equidade, escopo, espírito, estímulo, exigir, explicitar, expressivo, focar, fomentar, fragmentação, historicamente, histórico-cultural, hoje, horizonte, humanidade, ilustrar, implicar, inter-relação, interdisciplinaridade, lado, logo, metodológico, moderno, notar, novidade, operar, ora, patamar, pedagogia, permear, postar, pp, preciso, preconizar, prova, qualificar, que, renovar, salientar, singularidade, sistematização, sumo, supor, tarefa, teoria, termos, teórico, teórico-metodológico, transdisciplinaridade, tópico, um, urgente, ver
\end{textsample}
